\clearpage
\SourceLeaf{34}

\begin{center}\large BEWEIS DER IRRATIONALITÄT DER TANGENTEN RATIONALER BÖGEN IN EINER NEUEN GESTALT.\end{center}

\noindent\hspace*{\parindent}Die unendlichen Reihen

\[
P = 1-\frac{1}{2} . \frac{mm}{nn}+\frac{1}{2.4} . \frac{1}{1.3} . \frac{m^4}{n^4}-\frac{1}{2.4.6} . \frac{1}{1.3.5} . \frac{m^6}{n^6}+\text{u.\,s.\,w.}
\]
\[
\begin{aligned}
P_1 &= \frac{m}{n}-\frac{1}{2} . \frac{1}{1.3} . \frac{m^2}{n^3}\\ &\quad+\frac{1}{2.4} . \frac{1}{1.3.5} . \frac{m^5}{n^5}-\frac{1}{2.4.6} . \frac{1}{1.3.5.7} . \frac{m^7}{n^7}+\text{u.\,s.\,w.}
\end{aligned}
\]
\[
\begin{aligned}
P_2 &= \frac{1}{1.3} . \frac{m^3}{n.n}-\frac{1}{2} . \frac{1}{1.3.5} . \frac{m^5}{n^4}\\ &\quad+\frac{1}{2.4} . \frac{1}{1.3.5.7} . \frac{m^7}{n^6}-\frac{1}{2.4.6} . \frac{1}{1.3.5.7.9} . \frac{m^9}{n^8}+\text{u.\,s.\,w.}
\end{aligned}
\]
\[
\begin{aligned}
P_3 &= \frac{1}{1.3.5} . \frac{m^5}{n^3}-\frac{1}{2} . \frac{1}{1.3.5.7} . \frac{m^7}{n^5}\\ &\quad+\frac{1}{2.4} . \frac{1}{1.3.5.7.9} . \frac{m^9}{n^7}-\frac{1}{2.4.6} . \frac{1}{1.3.5.7.9.11} . \frac{m^{11}}{n^9}+\text{u.\,s.\,w.}
\end{aligned}
\]

\noindent und allgemein, von $\theta = 1$ an,

\[
\begin{aligned}
P_\theta = {}& \frac{1}{1.3.5\ldots(2\theta-1)} . \frac{m^{2\theta-1}}{n^\theta}
-\frac{1}{2} . \frac{1}{1.3.5\ldots(2\theta+1)} . \frac{m^{2\theta+1}}{n^{\theta+2}}
\\ &\,+\frac{1}{2.4} . \frac{1}{1.3.5\ldots(2\theta+3)} . \frac{m^{2\theta+3}}{n^{\theta+4}}
-\frac{1}{2.4.6} . \frac{1}{1.3.5\ldots(2\theta+5)} . \frac{m^{2\theta+5}}{n^{\theta+6}}+\text{u.\,s.\,w.},
\end{aligned}
\]

\noindent wo $m$, $n$ zunächst beliebige gegebene positive Grössen bedeuten, sind
offenbar alle, wenn nicht vom Anfange an, doch nach hinlänglich weiter
Fortsetzung convergent, und zwar mehr als jede fallende geometrische
Progression. Es hat daher jede einen endlichen bestimmten Zahlenwerth.

Ist schon das zweite Glied von $P_\theta$ kleiner als das erste, oder auch
nur diesem gleich, so wird von selbst das dritte kleiner als das zweite, das
vierte kleiner als das dritte u.\,s.\,w. In diesem Falle, wozu nur gefordert
wird, dass $2\theta+1$ nicht kleiner ist als $\dfrac{mm}{2nn}$, wird der
Werth von $P_\theta$ jedenfalls positiv, kleiner als das erste Glied, und
grösser als die Summe der beiden ersten sein.
